\section{结论}
\label{sec:conclusion}

本文从一个结构性局限出发：ZSAD 中的通用异常 prompt 具备良好的跨类别泛化，却只能在类别无关的宽泛语义上刻画“异常”，使异常侧语义过宽，导致规则纹理、花纹等显著偏离正常模板的正常结构被误响应为异常。我们指出，收紧异常侧所需的位置与形态信息无法由通用 prompt 预先写定，只能在测试阶段逐图像地从图像自身获取；并以 CLIP patch 特征的高频分量作为无标签条件下的可计算依据，在测试时逐图收紧异常侧文本表示，使其对齐当前图像的实际异常。方法冻结 CLIP、无训练、无反向传播，在五个基准上于像素级定位取得一致提升，且不在正常图像上引入额外误报。

\paragraph{局限。}
高频依据在极端低对比或运动模糊图像上信息量下降，此时收紧信号偏弱；本文亦未探索多层小波与跨尺度证据的联合。这些留待后续工作。（本文所有数值为预期目标值，最终结论以真实实验回填后为准。）
